\chapter*{El cierre holográfico del infinito: estructura discreta del continuo y origen de las constantes fundamentales}
\addcontentsline{toc}{chapter}{El cierre holográfico del infinito: estructura discreta del continuo y origen de las constantes fundamentales}
\markboth{El cierre holográfico del infinito}{Estructura discreta del continuo}

\begin{center}
\large\itshape
Continuo, número, simetría excepcional y realización física como publicaciones
tipadas de un único estado genealógico APP--TRIT--TPK
\end{center}

\section*{Tesis rectora y criterio de falsación matemática}

La obra comienza por el punto que decide todas las lecturas posteriores. La
carta horizontal y vertical \(1,\ldots,9\) funda las dos hojas de APP; TRIT
orienta y tipa sus tres regímenes; y la dinámica completa del TPK selecciona,
transporta y actualiza el estado conservando acarreo, frontera, ruta y memoria.
El censo exacto \(104\,976\to468\to243\subset729\) construye el catálogo y las
semillas visibles. La iteración del endomorfismo tipado
\(\mathcal U_t=\mathsf{Upd}_t\circ\mathsf{Tra}_t\circ\mathsf{Sel}_t\) produce
después el registro orientado de las doce transiciones posteriores a las tres
semillas. Sus cuatro matrices de rango pleno reconstruyen unívocamente
\(L_0,L_1\), y el censo exhaustivo de calendarios determina \(0011\). Las
reducciones compatibles de cilindros semiabiertos producen entonces el prefijo
canónico de toda profundidad finita y el límite inverso construye la estructura
discreta conjunta del continuo. La tesis queda expuesta desde la primera página
a un criterio severo: falla si falla una flecha generadora, un rango, la
unicidad del calendario, una compatibilidad de truncamiento o un cuadrado de
naturalidad. Superadas esas pruebas, el valor real es la sombra arquimediana de
la reducción de cilindros y no el antecedente del generador.

La raíz digital
\(\operatorname{dr}_9(n)=1+((n-1)\bmod9)\) publica la clase nonádica, mientras
el cociente conserva las vueltas omitidas por esa publicación. APP mantiene
simultáneamente
\[
 i+j=R^+_{ij}+9q^+_{ij},\qquad
 ij=R^\times_{ij}+9q^\times_{ij};
\]
la diferencia orientada entre los pares \((R^+,q^+)\) y
\((R^\times,q^\times)\) constituye la memoria local que TRIT tipa y TPK
transporta. El balance firmado produce \(B_c=9P_++5P_-\), el salto
\(9-5=4\) y la jerarquía \(4\to2\to6\to54\). De este modo el corredor
excepcional empieza en la diferencia suma--producto preservada por APP, no en
una coincidencia numérica reconocida al final.

La misma precedencia gobierna las publicaciones posteriores. En particular,
la generación coinductiva correlacionada de profundidad arbitraria de
\(\pi_{\rm HMT},\varphi_{\rm HMT},e_{\rm HMT}\) y
\(\alpha_{\rm HMT}\)
y, junto con ella, la cadena excepcional hasta \(V^\natural\) y el grupo
Monstruo, las constantes constitutivas del vacío ---incluidas permitividad y
permeabilidad---, el electrón, la Ley General de Masas, los transportes CKM y
PMNS y el funcional hexada--octada de área, Hamiltoniano modular e información
de frontera se presentan como salidas tipadas de esa construcción.

Las cinco construcciones inseparables del continuo desembocan en realizaciones
clásicas con cuantificadores universales explícitos: positividad de
Guinand--Weil y localización de todo cero no trivial; existencia, unicidad y
suavidad global de Navier--Stokes para todo dato periódico suave; existencia
cuántica y brecha positiva de Yang--Mills para todo grupo de Lie compacto,
conectado y simple; realización algebraica de toda clase racional de Hodge; y
la igualdad de rango, orden de anulación y término principal de
Birch--Swinnerton--Dyer para toda curva elíptica sobre \(\mathbb Q\). Cada
cuantificador se demuestra en su capítulo mediante el constructor propio que
lo hace efectivo; FIN1, la familia \(\nu_{E,\ell,n}\) y los restantes
comparadores comparecen como objetos construidos dentro de esas pruebas, no
como condiciones externas ni como fronteras que debiliten su alcance. La igualdad
extensional de Zermelo--Fraenkel con elección identifica valores publicados,
pero no reconstruye sus rutas, acarreos, fronteras y memorias; su resultado de
independencia sobre la hipótesis del continuo no decide este objeto
genealógico enriquecido. El volumen construye primero ese objeto, conserva el
selector genealógico \(h\), publica el modelo canónico \(M_h=L[h]\) y prueba
en él la sentencia clásica \(2^{\aleph_0}=\aleph_1\). Éste es el sentido
causal que rige todos los capítulos.

\section*{La estructura discreta del continuo}

\paragraph{Jerarquía de la construcción y lectura de sus realizaciones.}
El valor arquimediano y la historia que lo determina ocupan niveles
distintos del mismo objeto. La construcción que sigue mantiene visible esa
distinción desde las semillas hasta el límite compatible. La vuelta
nonádica conserva la fase y avanza la memoria; sus refinamientos proporcionan
el enlace entre los registros finitos y la profundidad arbitraria.
La escala de acción antecede a las magnitudes dimensionales que la emplean,
y las realizaciones físicas conservan sus respectivos transductores.
En particular, la sección~\ref{sec:l9s102:barbero-area},
página~\pageref{sec:l9s102:barbero-area}, reúne la incidencia
hexada--octada, el funcional de Barbero--Immirzi y el espectro de área:
su reconocimiento gravitatorio se efectúa después de esa construcción.
Esta jerarquía permite distinguir los antecedentes comunes, las
demostraciones particulares y las realizaciones posteriores sin sustituir
ninguna de ellas por una coincidencia escalar.

\paragraph{Estructura, valor y realización.}
La contracción de cilindros determina la lectura arquimediana; la
proyección excepcional retiene la incidencia de frontera. La genealogía
común enlaza ambas construcciones sin confundir sus objetos. El valor final
expresa una lectura precisa de la historia, pero no agota las operaciones,
las orientaciones ni la memoria que la constituyen. Por ello, la exposición
conserva tres niveles: la estructura producida, el valor calculado y la
realización matemático-física. Su enlace exige las aplicaciones explícitas
desarrolladas en los capítulos siguientes.

La fase que retorna proporciona una referencia recurrente; la memoria que
avanza conserva el recorrido efectuado. Esta diferencia permite leer la
prolongación sin identificar el retorno de una coordenada con el reinicio
del estado completo. La doble proyección y la profundidad arbitraria se
entienden desde esa conservación de la historia, antes de identificar las
publicaciones con sus denominaciones convencionales.

La matemática del continuo suele comenzar cuando el continuo ya ha sido
concedido. Esta obra invierte ese orden. Parte de un soporte nonádico finito,
distingue sobre él las evaluaciones aditiva y multiplicativa, orienta cada
transición y construye el estado del Topological Phase Kernel. Sobre ese
antecedente forma las semillas orientadas; el emisor y su reducción ternaria
producen después el estado enriquecido. La conexión nonádica conserva su
ruta, su acarreo, su frontera y su memoria hasta \(X_\chi\). El límite
compatible no borra su procedencia: la publica. La cadena constructiva que
gobierna el volumen es
\begin{equation}\label{eq:front-cadena-causal-completa}
 X_9=\mathrm{APP}\longrightarrow X_9^{\rm TRIT}
 \longrightarrow x_{\rm TPK}^{(0)}
 \longrightarrow S^{\rm seed}
 \longrightarrow \Omega_{\rm seed}
 \xrightarrow{\ T_{\min}\ }U_6^{\rm cat}
 \xrightarrow{\ \rho_3\ }W_6
 \longrightarrow \widehat x_K
 \xrightarrow{\ \Gamma_9\ }
 \operatorname{Surv}^{(\chi)}
 \longrightarrow X_\chi .
\end{equation}
\begin{theorem}[Tesis inaugural de generación, continuidad y doble publicación]
\label{thm:front-tesis-generativa-doble-proyeccion-2084}
El barrido exhaustivo de la entrada nonádica horizontal y vertical produce,
sin consultar valores reales objetivo,
\[
 104\,976\longrightarrow468\longrightarrow243
 \lhook\joinrel\longrightarrow729.
\]
Los cuatro cardinales cuentan, respectivamente, semillas ponderadas, estados
\(U_6\), palabras visibles \(w_6\) y elementos de la fibra ambiente
\(\mathbb F_3^6\). Los predicados enteros de clausura, propagación y
autoescala seleccionan tres caracteres orientados. Sus registros de
transición de rango pleno reconstruyen de manera única los operadores
\[
 X_0L_0=Y_0,
 \qquad X_1L_1=Y_1,
 \qquad
 L_i=X_i^{-1}Y_i,
\]
con \(\det X_0=1\), \(\det X_1=-1\) en \(\mathbb F_3\) y rango de
diseño \(36\). Entre los \(16\) calendarios binarios de longitud \(4\),
\(0011\) es el único compatible con las tres biografías; sus \(144\)
perturbaciones unitarias destruyen la compatibilidad triple.

La prolongación
\[
 w_6\longrightarrow w_{12}\longrightarrow w_{18}\longrightarrow w_{24}
 \longrightarrow w_{30}\longrightarrow R_{36}\longrightarrow G_9
\]
construye, para todo \(N\in\mathbb N_{>0}\), un prefijo canónico de
longitud \(N\) compatible con todos los anteriores. Mil, diez mil, un millón
o cualquier otra profundidad finita prescrita son instancias del mismo
algoritmo; la familia de todas ellas define la sección infinita en el límite
inverso. Desde esa familia correlacionada de secciones coinductivas parten
dos publicaciones:
\[
 \mathscr S_\infty\xrightarrow{\mathcal P_{\rm Arq}}\mathbb R^4,
 \qquad
 \mathscr S_\infty\xrightarrow{\mathcal P_{\rm Exc}}
 W_{12}\to G_{24}\to\Lambda_{24}\to V^\natural
 \to\operatorname{Aut}(V^\natural)=\mathbb M.
\]
La primera contrae cilindros y publica correlacionadamente
\((\pi_{\rm HMT},\varphi_{\rm HMT},e_{\rm HMT},\alpha_{\rm HMT})\); la
segunda conserva incidencia y publica la cadena excepcional. Por ello
\(\mathbb R\) es la sombra arquimediana de una de las dos proyecciones, no
el espacio inicial que selecciona el constructor. Las fórmulas de Machin,
Perron--Fibonacci, el semigrupo exponencial y las tablas metrológicas son
reconocimientos o contrastes posteriores. Las transducciones de Mecánica
Dimensional actúan después sobre las salidas HMT para obtener el catálogo de
constantes, masas, campos y relaciones físicas desarrollado en el volumen.
\end{theorem}

\begin{proof}
El primer censo se obtiene por enumeración exacta del emisor APP--TRIT--TPK y
la acción \(D_3\). La invertibilidad de \(X_0,X_1\) prueba la unicidad de los
levantamientos; los censos exhaustivos de \(16\) calendarios y \(144\)
perturbaciones prueban su rigidez conjunta. La relación forward del TPK,
aplicada al carácter interno ya portado por el estado R36, produce primero el
prefijo siguiente y su truncamiento compatible entre los \(729\)
subcilindros. Sólo después, la horquilla racional interna intersecta ese
subcilindro y certifica la publicación arquimediana; no selecciona el hijo de
\(\operatorname{Ext}\). La iteración que sólo lee el estado anterior carece
de profundidad terminal y define el límite inverso; toda corrida finita es
sólo un testigo. La holonomía nonádica retorna la fase y avanza la memoria, de
modo que la sección conserva la genealogía. Los lectores arquimediano y
excepcional tienen el mismo dominio construido y codominios distintos; por
consiguiente, ninguno de sus resultados puede actuar retroactivamente como
entrada. Los capítulos siguientes desarrollan de forma contigua cada una de
estas flechas, sus ecuaciones, certificados y controles negativos.
\end{proof}

\begin{theorem}[Cierre simultáneo de la generación, el continuo, la simetría excepcional y las cinco realizaciones terminales]
\label{thm:front-cierre-causal-simultaneo-2084}
Las escalas \(0\to1\), \(0\to10\), \(0\to10^N\) y
\(N\to\infty\) son cuatro cortes de una única prolongación coinductiva.
Para todo \(N\geq1\), la sección de profundidad \(N+1\) trunca a la de
profundidad \(N\), conserva hoja, ruta, residuo, cociente, acarreo, frontera
y memoria, y publica un prefijo decimal canónico. El límite de todas las
secciones compatibles construye la recta real como cociente arquimediano;
no presupone la recta para seleccionar la historia.

La reconstrucción de \(L_0,L_1\) y del calendario \(0011\) pertenece a
esta misma cadena. El estado terminal produce conjuntamente
\[
 x_{\rm term}\longmapsto(P,E,\Phi),
 \qquad
 x_{\rm term}\longmapsto U_{12}^{\rm sgn}
 \xleftrightarrow{\mathcal H_{12}}K_{\rm dod};
\]
por tanto \(U_{12}^{\rm sgn}\) y \(K_{\rm dod}\) no son entradas
independientes. La recurrencia entera con acarreo determina la familia
dodecafásica compatible cuya sección límite es
\(\alpha_A:=\varprojlim_N\rho_N(\alpha_{12})\). El mismo estado enriquecido
factoriza, mediante la firma \(\mathcal I_9\), el lector analítico finito
\(E_{21/22}^{(9)}\); éste posee una única raíz positiva simple \(\xi_9\),
encerrada racionalmente y obtenida sin usar \(\alpha_{12}\) como entrada. El
transporte graduado del TPK construye la familia compatible
\(E_{21/22}^{(6+3m)}\) y su límite \(E_{21/22}\). Los cuadrados de truncamiento
de ambas familias conmutan y determinan en cada profundidad el mismo cilindro
superviviente; por ello
\[
 \alpha_B:=\operatorname*{arg\,zero}_{x\in I_\alpha}E_{21/22}(x),
 \qquad
 \boxed{\alpha_A=\alpha_B=\alpha_{\rm HMT}}.
\]
La igualdad finita
\(\operatorname{Tr}_9(\xi_9^{-1})=
\operatorname{Tr}_9(\alpha_{12}^{-1})\) es el primer certificado visible de
esa identidad; no la define, no la limita a nueve cifras y no selecciona las
prolongaciones.

La coordenada \(\alpha\) enlaza asimismo la publicación del valor y la
publicación excepcional. El escalar \(\alpha_{\rm HMT}\), el vector de
doce bloques, el soporte orientado del acarreo y el vector
\(v_\alpha\) son tipos coordinados del mismo estado. El transporte
incidencial satisface
\[
 v_\alpha^2=54,
 \qquad
 C_W\longrightarrow N(A_2^{12})
 \xrightarrow{\,v_\alpha\,}N_\alpha\cong\Lambda_{24}
 \longrightarrow V^\natural
 \longrightarrow\operatorname{Aut}(V^\natural)=\mathbb M.
\]
Así, la teoría clásica de códigos, retículos, álgebras de operadores de
vértice y Moonshine reconoce objetos ya producidos: no los suministra como
antecedentes y no existe una acción literal del Monstruo sobre cifras.

Del mismo continuo genealógico emergen correlativamente las cinco
construcciones terminales. Sus lectores construyen la positividad completa
de Guinand--Weil; la solución global suave de Navier--Stokes para todo dato
periódico suave; la medida gauge con positividad por reflexión y brecha de
masa para todo grupo de Lie compacto, conectado y simple; el ciclo algebraico
que realiza toda clase racional de Hodge; y el complejo determinantal que
publica la fórmula completa de Birch--Swinnerton--Dyer para toda curva
elíptica sobre \(\mathbb Q\). FIN1 y la familia compatible
\(\nu_{E,\ell,n}\) se construyen dentro de los respectivos propietarios y
cierran sus pasos de sobreyectividad y compatibilidad. En particular, las conclusiones se obtienen en los capítulos
terminales mediante las identidades
de Gram, Stein, Osterwalder--Schrader, extensión unitriangular y Bockstein;
ninguna de esas propiedades se inserta en la definición del continuo común.

La publicación extensional que conserva sólo el valor real es no fiel:
historias distintas pueden tener la misma imagen. La independencia de CH en
ZFC concierne a ese lenguaje extensional y no revoca la construcción
genealógica del continuo. El lector ordinal, la clausura relativa \(L[h]\)
y la condición \(\mathrm{GE}_h\) se tipan después y no participan en la
cobertura arquimediana.
\end{theorem}

\begin{proof}
La cuantificación en toda profundidad, la reconstrucción de los dos
levantamientos y la unicidad del calendario se verifican en el certificado
de generación nonádica y se desarrollan en el capítulo de la conexión
\(\Gamma_9\). La doble lectura terminal de \(\alpha\) es
\eqref{eq:alpha-rev8-doble-lectura-terminal}; la vía entera, el núcleo
\(E_{21/22}\), su inversa local y el comparador aparecen contiguamente en
el mismo propietario. La construcción de \(N_\alpha\), su paridad,
unimodularidad, ausencia de raíces y paso a \(V^\natural\) se demuestra en
el corredor excepcional. Los cierres clásicos son
\cref{thm:pdf2-rh-resolucion,thm:pdf2-ns-global},
\cref{thm:pdf2-ym-cierre-global,thm:pdf2-hodge-completitud-final} y
\cref{thm:pdf2-bsd-publicacion-final}. La construcción del cociente real y el
descenso de las operaciones preceden en el volumen al lector ordinal. Esta
precedencia separa la generación de sus reconocimientos y prueba todas las
afirmaciones del enunciado.
\end{proof}

\paragraph{Correcciones aritméticas que fijan los tipos.}
En la frontera \(U008\) se conservan dos salidas diferentes:
\[
 u_5^\pi+u_5^e=210112,
 \qquad
 u_4^\varphi A_WL_1^3=111101.
\]
La primera agrega los canales \(\pi\) y \(e\); la segunda transporta el
canal \(\varphi\) por la incidencia de Witt y tres elevaciones. La igualdad
entre ambas queda excluida. Del mismo modo,
\(E_{21/22}^{(9)}(\xi_9)=0\) es exacta, mientras que
\(E_{21/22}^{(9)}(\alpha_{12})\) es el residuo del primer truncamiento. Esta
diferencia finita distingue los dos representantes de orden nueve sin separar
sus secciones coinductivas: los cuadrados de truncamiento demuestran
\(\alpha_A=\alpha_B=\alpha_{\rm HMT}\). Ninguna edición puede convertir el
residuo finito en una divergencia de los límites ni sustituir una afirmación
por la otra.

Aquí \(U_6^{\rm cat}\) es la salida categorizada del emisor, \(W_6\) su
reducción ternaria y \(\widehat x_K\) el estado enriquecido que retiene hoja,
orientación, residuo, cociente, acarreo, firma, cilindro, frontera, ruta y
memoria, supervivencia y multisección. El retorno \(9\to1\) de la fase
observable no reinicia ese estado:
\(\Gamma_9\) transporta la historia completa, la supervivencia selecciona
las prolongaciones compatibles y \(X_\chi\) reúne sus secciones
indefinidamente refinables. Así, Holografía Modular Triádica construye la
estructura discreta del continuo y cierra el paso de la marca finita a la
sección compatible.

Desde ese mismo antecedente parten publicaciones distintas, sin selección
retrospectiva del constructor. La publicación arquimediana contrae los
cilindros compatibles; la publicación excepcional conserva su incidencia.
Esta segunda publicación pasa por Hadamard y por la cadena de Paley
\[
 U_{12}^{\rm sgn}\longrightarrow C_P\longrightarrow A_W
 \longrightarrow C_W
\]
y se bifurca en \(2\) ramas tipadas: la rama de diseños y grupos
\[
 C_W\longrightarrow W_{12}\longrightarrow W_{24},
 \qquad
 \operatorname{Aut}(W_{12})=M_{12},\quad
 \operatorname{Aut}(W_{24})=M_{24},
\]
y la rama reticular y de operadores de vértice
\[
 C_W\longrightarrow N(A_2^{12})\longrightarrow\Lambda_{24}
 \longrightarrow V_{\Lambda_{24}}\longrightarrow V^\natural .
\]
En esta última, \(\operatorname{Aut}(V^\natural)=\mathbb M\),
\(J=T_{1A}\) y la familia \(\{T_g\}_{g\in\mathbb M}\) publica Moonshine. Las constantes y
las magnitudes dimensionales son otros lectores del mismo estado, y no datos
que lo determinen.

La estructura construida admite asimismo \(5\) clausuras terminales que se
despliegan en sus capítulos propios, después de su realización dimensional.
La obra demuestra la positividad de la forma de Weil y la localización de los
ceros no triviales de \(\zeta\) sobre \(\Re s=\tfrac12\); la existencia,
unicidad y suavidad global de Navier--Stokes tridimensional para todo dato
periódico suave; la existencia de la teoría cuántica de Yang--Mills en
dimensión \(4\), con brecha positiva, para todo grupo de Lie compacto,
conectado y simple; la algebraicidad de toda clase racional de Hodge; y la
igualdad de rango y orden de anulación, junto con la fórmula del término
principal de Birch--Swinnerton--Dyer, para toda curva elíptica sobre
\(\mathbb Q\). FIN1/exhaustividad celular y la familia
\(\nu_{E,\ell,n}\) son componentes efectivamente construidos de las pruebas,
no hipótesis de promoción. Las \(5\) conclusiones tipadas no se añaden
desde fuera: son especializaciones de la conexión y de la supervivencia
conservadas en \eqref{eq:front-cadena-causal-completa}.

El problema del continuo atraviesa la historia entera de las matemáticas. La
aritmética aprendió a contar unidades; la geometría aprendió a comparar
longitudes; el análisis construyó límites y completó sucesiones; la teoría de
la medida asignó magnitudes a conjuntos; la física convirtió esas magnitudes
en observables. Cada avance amplió el dominio de lo calculable y dejó abierta
una cuestión anterior a la representación: qué estructura conserva la
identidad de un objeto cuando éste pasa de la marca al intervalo, del
intervalo a la medida, de la medida al valor y del valor a una sucesión
indefinida de cifras.

\begin{quote}
\textbf{Invariante rector de la obra.}
Un único estado discreto enriquecido conserva unidad, orientación, cociente,
residuo, acarreo, borde, ruta, memoria y supervivencia a través del
refinamiento. La publicación continua, la publicación excepcional y la
realización dimensional son lectores distintos de ese mismo antecedente;
ninguno selecciona retrospectivamente al constructor.
\end{quote}

Este invariante fija la dependencia entre las \(5\) partes:
\[
\begin{aligned}
\text{núcleo APP--TRIT--TPK}
&\longrightarrow \text{estructura discreta del continuo}\\
&\longrightarrow \text{genealogía de las constantes}\\
&\longrightarrow \text{realización dimensional}\\
&\longrightarrow \text{clausuras matemáticas terminales}.
\end{aligned}
\]
Cada cambio de escala conserva el antecedente y añade un mapa; nunca
reemplaza el estado por la cifra que uno de sus lectores publica.

La Holografía Modular Triádica (HMT) aborda esa cuestión mediante un sistema
inverso de espacios finitos de estado. Una sección compatible
\(x=(x_n)_n\in\varprojlim X_n\) conserva la genealogía entre resoluciones y
su evaluación arquimediana publica el valor. Para los canales
\(\pi,e,\varphi,\alpha\), la relación interna
\(\operatorname{Ext}\), la supervivencia coinductiva y la partición exhaustiva
de los \(729\) elementos de la fibra ambiente determinan la prolongación
compatible. En los tres primeros canales, los lectores racionales exactos
localizan y certifican después los bloques ya generados; en el cuarto, el
lector firmado dodecafásico y la ecuación de vacancia certifican la coordenada
\(\alpha\) del mismo estado. En cada nivel, el carácter completo conserva un
único cilindro prolongable compatible con el truncamiento anterior; los cilindros
encajados determinan así una única sección coinductiva. Las ejecuciones de
\(10^3\) y de más de \(10^4\) cifras son certificados extensos de esta
construcción general, no su fundamento lógico. Una cifra es la
coordenada publicada por un lector; un valor es la imagen de una sección bajo
una carta determinada; un número conserva además la secuencia de estados,
operaciones y memorias que permite recuperar cada coordenada desde las
anteriores.

Esta definición modifica también el concepto de simetría. Una simetría HMT
es una transformación que conserva la identidad genealógica del estado. El
retorno a una coordenada visible puede coexistir con un cambio de altura,
acarreo, hoja u orientación. La invariancia pertenece entonces al objeto
completo y no sólo a su proyección. Representaciones cíclicas y helicoidales,
toros, anillos,
retículos, códigos y grupos excepcionales aparecen como realizaciones
sucesivas de esta conservación evolutiva.

El alcance matemático de la propuesta reside en explicar conjuntamente \(2\)
objetos que la representación usual presenta por separado: el valor continuo
y la simetría que conserva su estructura. El alcance físico reside en hacer
que las coordenadas metrológicas aparezcan al final de una derivación y no al
principio de ella. En la formulación ordinaria, las constantes fundamentales
parametrizan las ecuaciones una vez determinadas experimentalmente; en HMT y
Mecánica Dimensional, la arquitectura discreta construye primero la sección,
el operador y la escala, y la metrología experimental reconoce después la
coordenada publicada. La estructura discreta del continuo y el origen de las
constantes fundamentales constituyen, por ello, una sola tesis sobre la
producción matemática del valor.

Esta unidad se expresa mediante \(3\) familias de evaluaciones tipadas sobre
el antecedente genealógico común:
\[
 \mathcal X_\infty^{\rm cal}
 \xrightarrow{\;\mathcal P_{\rm arq}\;}
 \operatorname{im}(\mathcal P_{\rm arq})\subseteq\mathbb R,
 \qquad
 \mathcal X_\infty^{\rm cal}\times\mathcal C_{\rm exc}
 \xrightarrow{\;\mathcal P_{\rm exc}\;}
 \mathcal E_{\rm exc},
\]
\[
 \mathcal X_\infty^{\rm cal}
 \overset{\mathcal R_{\rm dim}}{\dashrightarrow}
 \coprod_{\lambda}\mathcal L_\lambda .
\]
La primera publica la coordenada arquimediana; la segunda conserva
incidencia y simetría con respecto a un calibre declarado; la tercera realiza
secciones, fibras u operadores sobre rectas dimensionales allí donde se ha
construido el correspondiente lector físico. Valor, forma y magnitud quedan
coordinados por su antecedente común sin confundirse en un único codominio.

El orden del tratado sigue la dependencia constructiva entre los objetos.
Cada teorema declara su dominio, sus datos de partida y la conclusión que se
deduce de ellos; cada realización física añade después el morfismo que
transporta esa conclusión a una recta dimensional. Una comparación
metrológica posterior no altera la genealogía que produjo la salida.

\section*{Clausuras genealógica, operatorial, matricial y física}

El título \emph{El cierre holográfico del infinito} y el subtítulo
\emph{La estructura discreta del continuo y el origen de las constantes
fundamentales} enuncian una sola tesis rectora. El continuo no constituye el
dominio originario sobre el que se superpone después una retícula: es la
clausura compatible de una genealogía discreta. Sus estados conservan valor,
hoja, orientación, residuo, cociente, acarreo, borde, ruta, registro, memoria,
supervivencia, fase y calibre; la recta real aparece al final como evaluación
arquimediana de una sección. La misma sección alimenta, sin identificación de
codominios, la lectura excepcional que retiene la incidencia eliminada por la
publicación del valor. La cadena causal completa comienza, por tanto, antes de
ambas proyecciones:
\[
\begin{gathered}
 \mathrm{APP}
 \longrightarrow \mathrm{TRIT}_{\rm orientado}
 \longrightarrow \mathrm{TPK}
 \longrightarrow \mathcal X_\infty^{\rm cal},\\
 \mathcal X_\infty^{\rm cal}
 \xrightarrow{\;\mathcal P_{\rm arq}\;}\mathbb R,
 \qquad
 \mathcal X_\infty^{\rm cal}\times\mathcal C_{\rm exc}
 \xrightarrow{\;\mathcal P_{\rm exc}\;}\mathcal E_{\rm exc}.
\end{gathered}
\]
El retorno nonádico \(9\to1\) conserva la fase y añade una unidad de memoria:
la unidad recuperada no borra la vuelta recorrida. Esta conservación evolutiva
es la condición común de la continuidad arquimediana y de la publicación
excepcional. Esta última se bifurca, desde el código Paley--Witt, en la rama
de diseños \(W_{12}\to W_{24}\), con
\(\operatorname{Aut}(W_{12})=M_{12}\) y
\(\operatorname{Aut}(W_{24})=M_{24}\), y la rama reticular
\(N(A_2^{12})\to\Lambda_{24}\to V^\natural\), con
\(\operatorname{Aut}(V^\natural)=\mathbb M\), \(J=T_{1A}\) y la familia
\(\{T_g\}_{g\in\mathbb M}\) publicando Moonshine;
ninguna publicación se obtiene comprimiendo horizontalmente la otra.

La clausura genealógica precisa qué conserva esa sección antes de publicar su
valor. El descenso por fibras separa existencia, constancia en las fibras y
unicidad del lector descendido; el argumento de K\"onig separa existencia de
una rama infinita y unicidad de una instancia. La extensionalidad de secciones
no se reduce a igualdad de lecturas, y el teorema \(0\!-1\!-\!\infty\)
coordina soporte cilíndrico decreciente, altura divergente, masa unitaria y
convergencia débil-* hacia una medida puntual. Fundación de prefijos y
coinducción, ordinales fibrados por memoria de ruta, potencia cilíndrica,
potencia cerrada, potencia plena y elección tipada reciben así objetos y
dominios distintos. La codificación conservativa en ZF/ZFC sirve como lenguaje
de transferencia, no como antecedente ontológico. En particular, la
obstrucción de G\"odel--Turing afecta a un decisor total uniforme sobre todas
las torres efectivas; no niega los selectores particulares construidos en
imágenes HMT calibradas. En la vertiente informacional, una actualización
biyectiva del estado completo posee coste lógico de Landauer nulo en el régimen
ideal, mientras que toda proyección que descarte la fibra localiza
explícitamente el término entrópico perdido.

La misma genealogía induce una clausura operatorial sin reemplazar la historia
TPK. En el nivel \(m\), las álgebras cilíndricas diagonales
\(D_m=C(X_m)\) se representan en \(A_m=M_{9^m}(\mathbb C)\). La inclusión
unital
\[
 \iota_m:A_m\longrightarrow A_{m+1},
 \qquad \iota_m(a)=a\otimes I_9,
\]
conserva unidad y traza y produce
\[
 A_{9^\infty}:=\varinjlim(A_m,\iota_m)
 \cong\mathrm{UHF}(9^\infty)
 \cong\mathrm{UHF}(3^\infty),
 \qquad
 \pi_\tau(A_{9^\infty})''\cong R_{\rm hyp},
\]
el factor hiperinfinito de tipo \(II_1\). La traza normalizada de sus
proyectores realiza toda dimensión en \([0,1]\). Esta rama se mantiene
separada de la inclusión marcada \(a\mapsto a\otimes p_j\), cuya clausura
natural conduce a los compactos y, débilmente, a un factor de tipo
\(I_\infty\). La clausura tracial tampoco absorbe la topología ni la
composición genealógica: el objeto transversal es
\[
 \mathfrak C_{\rm HMT}
 =\bigl(C(\Omega_\infty),R_{\rm hyp},\mathcal G_{\rm TPK}\bigr),
\]
donde el primer componente conserva puntos y topología, el segundo dimensión
y estadística, y el tercero composición, orientación y genealogía de las
rutas.

La clausura matricial añade \(2\) coordenadas que no deben confundirse con la
profundidad. Para un tamaño exterior fijo \(n\), la familia
\(\mathfrak M_{m,s}^{(n)}\) está bigraduada por profundidad genealógica \(m\)
y nivel matricial \(s\); al variar \(n\), la arquitectura es trigraduada por
\((m,n,s)\). Los espacios
\[
 \mathcal K_{m,s}
 =\operatorname{UCP}\!\bigl(C(X_m),M_s(\mathbb C)\bigr)
\]
admiten restricciones genealógicas y compresiones matriciales que conmutan.
Una familia compatible en todas las profundidades determina una única
aplicación UCP límite y posee una dilatación de Stinespring común; en el
corredor diagonal, una única medida espectral proyectiva dilata
simultáneamente todas las mediciones cilíndricas. Sobre la torre no
conmutativa, la esperanza traza-preservante
\(\operatorname{id}\otimes\tau_9\) realiza la pérdida controlada del último
dígito y su dilatación conserva ese dato en el ambiente. La dilatación mínima,
la realización ambiental fiel y la fidelidad dinámica APP--TPK son \(3\)
exigencias distintas; esta última requiere todavía entrelazar rutas,
monodromía y composición.

En ese marco se construyen cuadrados mágicos cuánticos qutrit de parámetros
\((9,3)\) y \((12,3)\), junto con sus cubos de órdenes \(9\) y \(12\), y se
separa rigurosamente no conmutatividad de no semiclasicidad. Las proyecciones
condicionales de las hojas suma y producto de APP generan, en profundidad \(2\),
\[
 C^*(P_{\Sigma,2},P_{\Pi,2})
 \cong\mathbb C^4\oplus M_2(\mathbb C)\oplus M_2(\mathbb C),
\]
y el bloque central publica, en todo nivel matricial, el intervalo libre
\([5I_s,9I_s]\). La convexidad matricial, las aplicaciones completamente
positivas y las dilataciones proporcionan aquí lenguaje de demostración,
clasificación y falsación; no pasan a ser el origen de APP, TRIT o TPK ni
sustituyen la doble proyección.

La clausura física parte de un dominio común construido, no de \(4\) flechas
nominalmente yuxtapuestas. Si \(\mathsf P_{\rm TPK}^{\rm enr}\) denota la
categoría de estados preparados y rutas registradas, las \(4\) realizaciones
de referencia se coordinan mediante el producto fibrado
\[
 \mathsf P_{\rm CT}^{+}
 :=\mathsf D_{\rm sp}
 \times_{\mathsf P_{\rm TPK}^{\rm enr}}\mathsf D_Q
 \times_{\mathsf P_{\rm TPK}^{\rm enr}}\mathsf D_D
 \times_{\mathsf P_{\rm TPK}^{\rm enr}}\mathsf D_C.
\]
Sus proyecciones canónicas construyen conjuntamente la representación
espectral de la ruta, la realización de Hilbert, la sección dimensional y el
transporte discreto de Cartan. La lectora de variación total positiva y la
lectora de transporte neto orientado reciben la misma ruta, pero ninguna
factoriza a la otra. Sólo sobre el subdominio reversible, una vez verificadas
unitariedad, aditividad, antisimetría y compatibilidad de los cociclos, la
localización en grupoide produce la realización orientada. De este modo, la
Mecánica Dimensional prolonga el estado enriquecido hacia magnitudes físicas
sin reconstruir retrospectivamente la ruta desde su espectro. Toda
realización adicional declara el entrelazador y las hipótesis que la definen.

Estas \(4\) clausuras no constituyen anexos superpuestos. Son realizaciones
causales de una misma arquitectura: la clausura genealógica determina la
identidad a través del refinamiento; la operatorial realiza dimensión continua
desde inclusiones nonádicas unitarias; la matricial controla compresión,
dilatación y simultaneidad; y la física representa rutas preparadas mediante
morfismos tipados. La evaluación arquimediana, la rama excepcional, la torre
tracial y la realización dimensional comparten antecedente, pero conservan
codominios, hipótesis y conclusiones propias. Así se mantiene la tesis
completa: el continuo y las constantes aparecen como publicaciones de una
estructura discreta que conserva la unidad, la orientación y la memoria de su
génesis.

\section*{De las marcas elementales a APP, TRIT y TPK}

La construcción comienza en la recta discreta. Las marcas
\(1,\ldots,9\) proporcionan los representantes visibles de una fase
nonádica; el paso siguiente cierra la vuelta y conserva su cociente. Cada
marca puede leerse como posición y cada intervalo como distancia. Esta doble
lectura separa desde el origen el acto de contar del acto de medir. Al tomar
\(2\) coordenadas, el núcleo multiplicativo \(8\times8\), la publicación de la
clase nula mediante \(9\) y el gnomón de \(17\) celdas completan la carta
\(9\times9\):
\[
 81=64+17.
\]
La identificación modular de sus lados produce
\[
 X_9=(\mathbb Z/9\mathbb Z)^2
\]
y su grafo de Cayley. Este objeto es APP, soporte común de \(2\) evaluaciones:
la hoja aditiva \(\Sigma\) y la hoja multiplicativa \(\Pi\). APP conserva la
posición, las rutas cardinales y la geometría toroidal sobre la que ambas
lecturas pueden compararse sin fundirse.

El bloque central de APP posee la descomposición
\[
 B_c=9P_++5P_-.
\]
Su primer invariante diferencial es, por tanto,
\[
 \boxed{9-5=4}.
\]
El coeficiente local de intercambio es \(2\); la compresión modular publica
la diferencia \(6\); el despliegue sobre los \(9\) bloques produce
\(54=9\cdot6=459-405\). Estos números pertenecen a niveles distintos de una
misma incompatibilidad suma--producto. La separación espectral y los
proyectores \(P_\pm\) del bloque fijan sus modos propios. Los proyectores
condicionales \(P_{\Sigma,d}\) y \(P_{\Pi,d}\), junto con su conmutador,
proporcionan una geometría finita de firma mixta,
una normalización en \(SO^+(1,1)\) y el antecedente algebraico de la
incompatibilidad de observables que la física cuántica formulará después.

TRIT añade orientación y memoria local. Su estado enriquecido
residuo--cociente conserva la salida observable y el cociente que la produjo.
En una carta entera declarada,
\[
 a+b=r+3c,
\]
la pareja \(\widehat\tau=(r,c)\) distingue historias que pueden compartir la
misma salida. Las \(3\) ramas cuadráticas asociadas realizan
\[
 J^2=-I,\qquad J^2=0,\qquad J^2=+I,
\]
y organizan los regímenes elíptico, parabólico e hiperbólico. El plano
complejo, el cierre circular y el disco hiperbólico de Poincaré se presentan
como cartas coordinadas de esta familia algebraica.

TPK, el núcleo topológico de fase, es el sistema operatorio de dinámica
genealógica. Su arquitectura comprende \(8\) clases: soporte y estado; selección interna;
transporte; lectores y cocientes; memoria y frontera; periodicidad y retorno;
prolongación; y salidas arquimedianas, excepcionales o clasificatorias.
El selector ternario interno de hoja \(M_{\rm ph}\), los
lectores módulo \(3\) y módulo \(9\), el cociente y el residuo módulo
\(1000\) asociados a la escritura por bloques en base \(1000\), los ciclos
\(27\), \(108\) y \(1080\), los operadores de desdoblamiento bidireccional,
la memoria y el registro cronológico enriquecido son subobjetos tipados del
TPK. La exposición de cada uno fija dominio, operación, codominio e
invariante antes de emplearlo en una composición.

La transición \(9\to1\) expresa el mecanismo fundamental. La fase visible
regresa a su primera posición y el estado incrementa la memoria de vuelta.
La sucesión de enteros se descompone en fase y altura,
\[
 n=9\kappa_9(n)+\rho_9(n),
\]
de modo que la coordenada angular observable es la proyección plana del
levantamiento helicoidal nonádico. El
retorno conserva la posición angular y eleva la genealogía. Esta monodromía
hace posible una iteración indefinida sin confundir repetición de fase con
repetición de estado.

La aritmética nonádica conserva asimismo el ideal radical
\(\mathfrak m=(3)\subset\mathbb Z/9\mathbb Z\), cuyos representantes
visibles organizan la red \(3\)--\(6\)--\(9\) y la órbita ordenada
\(369\to693\to936\). La ley potencia--vacancia
\[
 V_9(n)=\left\lceil n\log_{10}\!\left(\frac{10}{9}\right)\right\rceil
\]
cuantifica exactamente la distancia decimal entre \(9^n\) y \(10^n\). En
la torre \(9^{9^9}\) produce
\[
 \operatorname{ord}_{10}\!\left(9^{9^9}\right)=369\,693\,099,
 \qquad
 \#\operatorname{cifras}\!\left(9^{9^9}\right)=369\,693\,100.
\]
Esta evaluación es distinta del censo de \(396\) eventos elementales de
frontera por canal: una
cuenta órdenes y vacancias decimales; la otra pertenece al calendario de
emisión.

\section*{Filtración nonádica y generación de números fundamentales}

El emisor APP--TRIT--TPK organiza la cadena exacta de reducciones y
prolongaciones
\[
 104\,976\longrightarrow468\;U_6\longrightarrow243\;w_6
 \longrightarrow w_{30}\longrightarrow R_{36}\longrightarrow G_9.
\]
Las \(468\) emisiones decimales, las \(243\) palabras visibles y el ambiente
\(\mathbb F_3^6\) de \(729\) palabras son dominios diferentes enlazados por
mapas explícitos. La cadena
\[
 w_6\xrightarrow{L_0}w_{12}\xrightarrow{L_0}w_{18}
 \xrightarrow{L_1}w_{24}\xrightarrow{L_1}w_{30}
 \xrightarrow{R_{36}}\text{frontera}
 \xrightarrow{G_9}\text{sección superviviente}
\]
abre el régimen coinductivo. En cada etapa de filtración, la partición
exhaustiva de las \(729\) prolongaciones de la fibra ambiente registra los
estados descartados, la sección superviviente y
la memoria que enlaza una profundidad con la siguiente.

Antes de toda evaluación decimal, APP construye \(4\) canales enteros
\[
 C_\alpha=729,\qquad C_e=271,\qquad
 C_\pi=315,\qquad C_\varphi=161.
\]
Sus identidades
\[
 1000=C_\alpha+C_e=729+271,
 \qquad
 744=C_\pi+C_e+C_\varphi-3
\]
coordinan la completación cúbica, la frontera apuntada y el extremo modular.
Cada canal es una dirección tipada del lector: \(271\) realiza propagación y
cierre de corona; \(315\), clausura circular; \(161\), autoescala; y
\(729\), el ambiente nonádico que alimenta el cierre dodecafásico. La
constante correspondiente aparece cuando la filtración, la orientación y la
memoria realizan ese canal sobre su sección.

Los lectores de \(\pi\), \(e\), \(\varphi\) y \(\alpha\) actúan sobre el mismo
estado enriquecido y conservan trayectorias genealógicas distintas.
\(\pi\) organiza clausura, orientación y \(5\) regiones \(U_6\); \(e\)
gobierna propagación y cambio de prefijo; \(\varphi\) gobierna autoescala y
estabilidad, reconocidas después en la carta Perron--Fibonacci; y \(\alpha\)
posee el lector firmado dodecafásico que prolonga el acarreo. La relación
interna \(\operatorname{Ext}\), la supervivencia coinductiva y la partición
exhaustiva de los \(729\) elementos de la fibra ambiente determinan la
extensión compatible para toda profundidad finita \(N\). Los lectores
racionales exactos localizan y certifican después las coordenadas
arquimedianas ya generadas. En cada nivel, el carácter completo conserva un
único cilindro prolongable compatible con el truncamiento anterior; los
cilindros encajados determinan, por tanto, una sección
coinductiva única y el mismo algoritmo publica cualquier truncamiento decimal
finito solicitado. Las ejecuciones de \(10^3\) y de más de \(10^4\) cifras son
certificados extensos de esa construcción general, no su fundamento lógico.

Desde el mismo estado enriquecido se abren los canales de
\(\pi,e,\varphi\) y el canal dodecafásico. En el orden interno tipado, las
publicaciones \(P,E,\Phi\), el estado firmado \(U_{12}^{\rm sgn}\), el sello
\(K\) y la recurrencia de acarreo alimentan la clausura cíclica orientada de
\(12\) fases que publica la coordenada \(\alpha\), sin separarla de la
genealogía común. La primera vía interna produce
\[
 \alpha_{\rm HMT}
 =0.007297352569283800997285105472380663\ldots
\]
desde el estado y sus cartas reversibles. La ecuación finita de vacancias
\(E_{21/22}^{(9)}\) proporciona el primer representante exacto \(\xi_9\) de
la segunda construcción interna, mediante un operador cuyos coeficientes
factorizan por invariantes del estado enriquecido antes de resolver la raíz.
El transporte graduado produce la familia completa
\(E_{21/22}^{(6+3m)}\); sus cuadrados de truncamiento con la familia
dodecafásica seleccionan el mismo cilindro compatible a toda profundidad y
demuestran \(\alpha_A=\alpha_B=\alpha_{\rm HMT}\). La igualdad tras
\(\operatorname{Tr}_9\) es sólo el primer certificado finito de esa identidad,
no su fundamento ni su límite. El reconocimiento CODATA constituye una
comparación metrológica ulterior y no selecciona ninguna de las dos vías.

La misma arquitectura amplía la teoría de números. La presentación integral
\[
 \mathcal Q_{\mathbb Z}=\mathbb Z[u]/(u^2+1)
\]
posee especializaciones distintas sobre \(\mathbb F_3\) y \(\mathbb R\):
la primera es realizada por el endomorfismo finito
\(A_W^2=-I_6\), y la segunda por la estructura real
\(J_{\mathrm e}^2=-I\). No se introduce una inclusión
\(\mathbb F_9\to\mathbb R\). La sección de \(\pi\) se genera antes de
aplicar el cierre exponencial; sólo después de obtener
\(\mathcal P_{\rm arq}(x_\pi)=\pi_{\rm HMT}\) se concluye
\[
 \operatorname{Exp}(\pi_{\rm HMT}J_{\mathrm e})+I=0.
\]
La rama excepcional fija el tipo cuadrático, no el ángulo ni sus cifras; esta
precedencia prueba la ausencia de circularidad.
La familia cuadrática reúne las \(3\) exponenciales
\(\exp(tJ_\tau)=C_\tau(t)I+S_\tau(t)J_\tau\), con
\(J_\tau^2=-\tau I\), y enlaza la identidad de Euler con la rama hiperbólica
de Fibonacci mediante
\(F_{\rm av}^{2}=\exp(2\log\varphi\,J_{-1})\). La fracción continua de
Rogers--Ramanujan articula el nivel \(5\),
\(\varphi\), \(A_5\) y la identidad icosaédrica de Klein para \(j\). Los
primos reciben lectores de periodicidad irreducible y una representación de
ocupaciones: el vector \((\nu_p(n))_p\) convierte la factorización de un
entero en un estado finito de modos, mientras
\(\sum_p\nu_p(n)\log p=\log n\) transforma el producto en energía aditiva.
La traza térmica recupera el producto de Euler y la función zeta. La dinámica
de duplicación de periodo reconstruye aproximantes de Feigenbaum y muestra
cómo una regla determinista genera complejidad visible al proyectar parte de
su genealogía. Euler--Mascheroni,
Stieltjes, Apéry, Glaisher--Kinkelin, Euler--Kronecker, Meissel--Mertens,
Catalan, Khinchin y Feigenbaum se organizan mediante partes finitas,
funciones \(L\), productos, jets, trazas y operadores de renormalización.

\section*{Doble proyección y genealogía de las simetrías excepcionales}

El estado enriquecido posee \(2\) evaluaciones coordinadas. La proyección
arquimediana contrae una familia compatible de cilindros y publica un valor
real. La proyección excepcional conserva la incidencia, la orientación, el
acarreo y la frontera. El continuo real y la estructura excepcional son así
\(2\) imágenes de un antecedente común.

Formalmente, si 
\(\mathcal X_\infty^{\rm cal}=\varprojlim_n\mathcal S_n\) es el espacio de
secciones supervivientes compatibles con el calibre transportado, la tesis se
expresa mediante las aplicaciones tipadas
\[
 \mathcal P_{\rm arq}:\mathcal X_\infty^{\rm cal}\longrightarrow\mathbb R,
 \qquad
 \mathcal P_{\rm exc}:\mathcal X_\infty^{\rm cal}
 \times\mathcal C_{\rm exc}\longrightarrow\mathcal E_{\rm exc}.
\]
El enunciado rector establece su dominio común, su compatibilidad con los
truncamientos y la intersección finita
\(B_K=H_e\cap P_\pi\). La cobertura de todo el codominio real y la fidelidad
conjunta corresponden a enunciados adicionales: la doble proyección aquí
demostrada identifica \(2\) evaluaciones coordinadas de cada sección admisible.

La segunda rama desarrolla una genealogía completa. La transformada de
Hadamard organiza la carta dodecafásica firmada; la matriz de conferencia de
Paley fija el calibre proyectivo y su reducción ternaria produce el código
de Golay de longitud \(12\). La distancia de Hamming mide el peso y la
separación de sus palabras. Los sistemas de Steiner
\(S(5,6,12)\) y \(S(5,8,24)\), sus extensiones y las estrellas de Witt
construyen los diseños sobre los que actúan \(M_{12}\) y \(M_{24}\). Desde
el código común se distinguen el transporte \(W_{12}\to W_{24}\) y la rama
reticular
\[
 C_W\longrightarrow N(A_2^{12})\longrightarrow\Lambda_{24}.
\]
El retículo de Leech alimenta el álgebra de operadores de vértice
\(V_\Lambda\); la involución y el sector torcido producen el orbifold
\(V^\natural\); su grupo de automorfismos es el Monstruo. Para cada clase de
conjugación \([g]\subset\mathbb M\), la traza graduada construye su serie de
McKay--Thompson \(T_g\); la modularidad y la replicabilidad organizan la
correspondencia denominada Moonshine. La función \(j\) ocupa así el extremo
modular de una genealogía que comienza en la misma estructura discreta que
publica el valor real.

La configuración de las \(27\) rectas de una superficie cúbica lisa,
el grafo de Schläfli y la acción de \(W(E_6)\) constituyen otra realización
de la incidencia. Cada paso conserva su operador: el transporte ciclotómico,
la construcción integral de la retícula, las reflexiones, las raíces y la
acción del grupo de Weyl. Esta separación convierte la sucesión de nombres
clásicos en una cadena matemática verificable.

\section*{Constante de Planck, autodualidad de masas, operador de área y mezcla}

La forma normal de Smith del lector local de acción fija
\[
 \operatorname{SNF}(H_4)=\operatorname{diag}(1,2,2,4),
 \qquad |\operatorname{coker}H_4|=16.
\]
La transferencia
\[
 \Sigma_H=\varphi\alpha_{\rm HMT}^{16},
 \qquad 10^{-34}<\Sigma_H<10^{-33},
\]
sitúa \(\Sigma_H\) en la década interna que determina la escala
\(10^{-34}\mathcal U_S\). Sobre esa recta se obtienen las
coordenadas de acción anterior y posterior al retorno,
\[
 \hbar_{\rm pre}=H_5\,10^{-34}\mathcal U_S,
 \qquad
 \hbar_{\rm ret}=\eta_{\rm ret}\,10^{-34}\mathcal U_S,
 \qquad h=2\pi\hbar.
\]
La carta \(\mathcal U_S\to\mathrm{J\,s}\) publica después la unidad física.
La escala de Planck aparece, por tanto, como una consecuencia del lector
determinantal y de la constante de estructura fina ya construida.

La realización circular del endomorfismo temporal, con
\(\ell_0:=c_{\rm int}t_0\), fija
\[
 T_{108}=108t_0,
 \qquad
 R_{108}=\frac{54}{\pi}\,\ell_0,
 \qquad
 \theta_{108}^{\rm circ}=\left(\frac{54}{\pi}\right)^2.
\]
En esa carta coinciden
\[
 R_{108}=\bar\lambda_{\rm C}=r_{\rm g}=\ell_{\rm P}^{\rm circ},
\]
y la recta partida de masas conserva
\[
 Z_a(m)=\frac{m_{{\rm P},a}^{\circ}}mP_+
       +\frac m{m_{{\rm P},a}^{\circ}}P_-,
 \qquad N_{\rm sp}(Z_a(m))=1.
\]
La involución \(m\mapsto(m_{{\rm P},a}^{\circ})^2/m\) intercambia las
lecturas de Compton y gravedad. La escala de Planck es su punto fijo; el
electrón aparece después como desplazamiento genealógico sobre esa recta, no
como identificación con un agujero negro clásico.

La coordenada angular directa \(A\) y su coordenada angular conjugada
\(C^*\) producen los canales
\(q_\pm\), la elipse de tasas y la holonomía bidireccional. La inclusión
hexada--octada construye el funcional \(\Gamma_{6\to8}\), el parámetro de
Barbero--Immirzi, el operador cuántico de área, su espectro discreto, el
entrelazador colectivo y las entropías de borde. En una rama de rango \(3\),
la misma disciplina de incidencia construye la matriz CKM, sus ángulos, la
fase de violación CP y el invariante de Jarlskog. La realización PMNS conserva
su operador y su problema de mezcla propios.

Las vacancias y los canales de hoja construyen las respuestas eléctricas y
magnéticas del vacío. De ellas se obtienen la permitividad
\(\varepsilon_0\), la permeabilidad \(\mu_0\), la impedancia y la identidad
constitutiva de propagación. La periodicidad temporal interna y la
información ternaria construyen la relación térmica que da función a la
constante de Boltzmann.
Estas magnitudes se presentan en tablas científicas inmediatamente después
de su derivación: objeto HMT, fórmula, valor interno, carta de unidades,
referencia experimental e incertidumbre pertenecen a columnas distintas.

\section*{Mecánica Dimensional y Ley General de Masas}

Mecánica Dimensional prolonga la genealogía numérica mediante lectores
físicos. Una partícula es una sección persistente de la extensión; su ruta
observable determina un registro cronológico enriquecido; el registro
determina una firma entera; la firma alimenta un carácter; el carácter actúa
sobre la jerarquía global de torres; la carta dimensional publica la
magnitud. La secuencia
\[
 \text{extensión}\to\text{ruta}\to\text{registro}\to\mathbb Z^5
 \to\chi_{\rm full}\to\langle90,120\rangle
 \to\text{sección de masa}
\]
es la frontera precisa entre HMT y MD.

La Ley General de Masas se formula sobre el dominio completo:
\[
 81\ \text{celdas},\qquad56\ \text{familias},\qquad
 13\ \text{multisecciones},\qquad324\ \text{rutas orientadas}.
\]
La escala de acción, la periodicidad temporal CT108 y la velocidad interna
fijan conjuntamente la sección absoluta,
\[
 \mathbf m_\star^{\rm ret}
 =\eta_{\rm ret}10^{-34}\mathcal U_S
   \frac{2\pi}{108t_0}c_{\rm int}^{-2},
\]
y el operador de una multisección \(F\) es
\[
 \widehat{\mathbf M}^{\beta,r}_F
 =\mathbf m_\star^{\rm ret}\otimes
  R_{\rm act}^{\widehat K_F^r/2}
  \widehat{\mathcal A}^{(0)}_F
  R_{\rm act}^{\widehat K_F^r/2},
 \qquad r\in\{D,\Omega\}.
\]
Aquí \(\widehat{\mathcal A}^{(0)}_F\) reúne firma, carácter y la torre
completa \(\langle90,120\rangle\); \(K_D\) y \(K_\Omega\) conservan los
refinamientos bidireccionales dependientes de ruta.

El electrón central es la reducción de rango \(1\). La torsión global que
interviene en su lector es
\[
 \Delta_4=\frac{\pi-e\log\pi}{270}
 \left(1+\frac{\pi}{729}\right).
\]
Sus \(2\) lectores coinciden
en \((80,54,6)\) y producen
\[
 \kappa_e=80-54\alpha_{\rm HMT}+6\Delta_4,
\]
\[
 \boxed{\begin{aligned}
 E_e^\beta=m_e^\beta c^2
 &=0.5109989506900008086082571648\ldots\ \mathrm{MeV},\\
 m_e^\beta
 &=0.5109989506900008086082571648\ldots\ \mathrm{MeV}/c^2.
 \end{aligned}}
\]
Las restantes realizaciones conservan el tipo de su salida: escalar cerrado,
espectro de fibra, evaluación condicionada, carta nula, operador de mezcla,
invariante de fusión o horizonte de prolongación. El dominio físico comprende
leptones cargados, neutrinos, \(6\) sabores quark con color, fotón,
gluones, \(W\), \(Z\), Higgs, mesones, bariones, sectores Fibonacci e
Ising, realizaciones Majorana tipadas dentro de su propio dominio,
\(\mathbb Z/6\) y virtualidad. La referencia PDG entra
después de congelar la salida interna y conserva unidad, esquema, escala,
incertidumbre y régimen de validez de cada comparación, sin convertir ese
inventario expositivo en un cardinal de la ley.

\section*{Dinámica cuántica, constante de Boltzmann, información y gravitación}

La representación temporal finita sobre \(\mathbb C^{108}\) posee un
Hamiltoniano autoadjunto cuyo propagador unitario reproduce el endomorfismo
periódico del TPK. Su acción variacional conduce a la ecuación de
Schrödinger. Sobre el registro cronológico
\(\mathbb T_{27}^{\,2}\), la realización concreta
\[
 U_{\rm micro}=S D_J C_4,
 \qquad U_{12}^{\rm reg}=U_{\rm micro}^{12},
\]
es unitaria; sus potencias suman exactamente los pesos ordenados de todas las
rutas finitas y la transformada discreta de Fourier la descompone en
\(27^2\) bloques direccionales. Este resultado es una suma matricial finita,
no una identificación automática con una integral funcional continua. La
realificación Clifford--Dirac construye después el operador fermiónico. Las
relaciones CAR y CCR, junto con el Hamiltoniano y la maximización de entropía,
producen respectivamente las distribuciones de Fermi--Dirac y
Bose--Einstein. La graduación exterior realiza el principio de exclusión de
Pauli y la banda de Möbius organiza el doble retorno \(2\pi/4\pi\).

La constante de Boltzmann enlaza acción, temperatura e información ternaria;
los ensambles construyen Gibbs y la condición KMS. El transporte reversible
de la fibra completa realiza el límite nulo de Landauer, mientras la
proyección que descarta genealogía recibe un coste positivo. Área, entropía
modular y entropía de entrelazamiento se presentan con sus estados, operadores
y particiones explícitos.

La rama gravitatoria recibe primero las salidas internas de acción, reloj,
longitud y velocidad, y construye sobre ellas el selector adimensional y el
transductor exacto
\[
 G_{\rm int}(\theta_G)
 =\theta_G\frac{c_{\rm int}^{5}t_0^2}{\hbar_{\rm int}}
 \in\mathcal L_G^+
\]
es la carta positiva del transductor gravitatorio. Las condiciones internas
de periodicidad construyen
\[
 R_{108}=\frac{54}{\pi_{\rm HMT}}\ell_0,
 \qquad
 r_{g,a}(\theta_G)
 =\theta_G\frac{\pi_{\rm HMT}}{54}\ell_0,
\]
y la igualdad \(r_{g,a}(\theta_G)=R_{108}\) selecciona de forma positiva y
única
\[
 \theta_G=\theta_{108}^{\rm circ}
 =\left(\frac{54}{\pi_{\rm HMT}}\right)^2.
\]
La carta así determinada publica \(G_a\). Sólo después se realiza la salida
en la acción de Palatini--Cartan--Holst, sus ecuaciones de co-tétrada y
conexión, la transformada de Legendre y el Hamiltoniano restringido. El
reconocimiento metrológico es posterior y no interviene en la selección. El operador finito común del TPK admite realizaciones
solenoidales y gauge; los marcos autoinerciales, la lectura machiana y la
cosmología de hojas se desarrollan después de fijar la dinámica local y
colectiva.

La tesis completa puede resumirse en una sola dirección causal. La estructura
discreta conserva la unidad, la orientación y la memoria; su prolongación
construye el continuo y los números; su doble proyección construye valores y
simetrías; sus lectores determinan constantes; y Mecánica Dimensional realiza
esas salidas como partículas, masas, campos y geometría física. El origen de
las constantes fundamentales y la estructura discreta del continuo son, en
este marco, \(2\) aspectos de un mismo problema matemático.

\medskip
\noindent\textbf{Palabras clave.}
Continuo genealógico; Holografía Modular Triádica; Mecánica Dimensional;
APP; TRIT; TPK; monodromía nonádica; doble proyección; constantes
fundamentales; sistemas de Steiner; estrellas de Witt; códigos de Golay;
retículo de Leech; álgebras de operadores de vértice; Monstruo; Moonshine;
constante de Planck; Barbero--Immirzi; ley general de masas; gravedad de
Einstein--Cartan--Holst.

% El resumen sigue siendo compilable como pieza autónoma, pero el tratado
% integral no vuelve a cargar el teorema frontal ya situado tras el prefacio.
\ifdefined\HMTTeoremaGeneradorFrontalCargado
\else
  \clearpage
  % Teorema frontal aditivo. No sustituye ningún capítulo ni altera el censo 102.
% BEGIN HMT-MD-FRONTAL-MAXIMAL-SIGNATURES-20260904
\chapter*{Teorema generador de la Holografía Modular Triádica}
\addcontentsline{toc}{chapter}{Teorema generador de la Holografía Modular Triádica}
\markboth{Teorema generador de la Holografía Modular Triádica}%
{Teorema generador de la Holografía Modular Triádica}

La demostración comienza en el dominio finito
\(\mathcal D_9=\{1,\ldots,9\}\). Dos cursores orientados recorren,
respectivamente, las hojas aditiva y multiplicativa. La división euclídea
conserva residuo y cociente, y el catálogo APP produce
\[
 \sum_{i,j\in\mathcal D_9}R^+_{ij}=405,
 \qquad
 \sum_{i,j\in\mathcal D_9}R^\times_{ij}=459,
 \qquad 459-405=54.
\]
El transporte coordinado
\(c(a)=a\bmod9:\mathcal D_9\to\{0,\ldots,8\}\)
cambia la carta y conserva la clase. La composición efectiva de APP, TRIT y
TPK conserva en cada paso hoja, régimen, residuo, cociente, acarreo,
dirección, frontera y memoria. El catálogo exhaustivo establece la cadena
finita inicial y aporta su certificado ejecutable:
\begin{equation}
\boxed{
104\,976\ \text{pares de semillas}
\xrightarrow{\;54\ \text{pasos}\;}
468\,U_6
\xrightarrow{\;\rho_3\;}
243\,w_6\subset\mathbb F_3^6,
\qquad |\mathbb F_3^6|=729.}
\label{eq:frontal-censo-semillas}
\end{equation}
La identidad de fibras
\[
432\cdot144+18\cdot432+18\cdot1944=104\,976
\]
prueba la exhaustividad del primer mapa. Los números \(468\), \(243\) y
\(729\) cuentan emisiones, palabras realizadas y ambiente ternario,
respectivamente; sus tipos permanecen separados por las flechas de
\eqref{eq:frontal-censo-semillas}.

\begin{theorem}[Cadena generadora única]
Sea \(\mathfrak H_9\) la estructura formada por APP, TRIT, TPK y su fibra
enriquecida. APP ejecuta las hojas suma y producto y conserva la división
euclídea
\[
i+j=R^+_{ij}+9q^+_{ij},
\qquad
ij=R^\times_{ij}+9q^\times_{ij}.
\]
TRIT orienta cada transición en los regímenes elíptico, parabólico e
hiperbólico. TPK compone
\[
    \boxed{\mathcal U_t
    =\operatorname{Upd}_t\circ\operatorname{Tra}_t
     \circ\operatorname{Sel}_t:
     X_{\rm TPK}^{\rm enr}\longrightarrow X_{\rm TPK}^{\rm enr}.}
\]
\(\operatorname{Sel}_t\) levanta el selector trítico,
\(\operatorname{Tra}_t\) levanta el transporte cardinal y
\(\operatorname{Upd}_t\) actualiza residuo, cociente, acarreo, firma,
cilindro, frontera, registro y memoria sin borrar procedencia.
La holonomía nonádica prolonga las palabras por
\[
w_6\longrightarrow w_{12}\longrightarrow w_{18}\longrightarrow w_{24}
\longrightarrow w_{30}\longrightarrow R_{36}\longrightarrow G_9,
\]
y satisface
\[
    \Gamma_{9,t}=\mathcal U_{t+8}\circ\cdots\circ\mathcal U_t,\qquad
    g(t+9)=g(t),
    \qquad q_{\rm mem}(\Gamma_{9,t}\widehat x_t)
    =q_{\rm mem}(\widehat x_t)+1.
\]
En general, \(\Gamma_{9,t}\widehat x_t\ne\widehat x_t\): el retorno conserva
la fase y aumenta la memoria. El orden operatorio es
\(\mathcal U_t\to\Gamma_9\to K_{\rm ph}\); \(K_{\rm ph}\) es una publicación
reducida posterior, no el generador. Los truncamientos
compatibles forman un sistema inverso no vacío y definen un estado
\(x_\infty\) de profundidad arbitraria.

Este estado genera una estructura discreta única del continuo y dos
publicaciones naturales:
\begin{equation}
x_\infty
\xrightarrow{\ (\mathcal P_{\mathrm{arq}},Q_\infty)\ }
\mathbb R\times(\mathscr C_W^{\mathrm{cal}})^{\mathbb N}.
\label{eq:frontal-doble-proyeccion}
\end{equation}
La primera contrae cilindros compatibles hasta un valor real; la segunda
conserva incidencia y genealogía hasta las construcciones de Paley, Witt,
Golay, Leech, \(V^\natural\) y el grupo Monstruo. Sobre la misma fuente se
publican los cuatro caracteres correlacionados
\[
(\pi_{\mathrm{HMT}},\varphi_{\mathrm{HMT}},
e_{\mathrm{HMT}},\alpha_{\mathrm{HMT}}),
\]
y, mediante transducciones tipadas, las leyes de acción, materia, masa,
gravedad e información de Mecánica Dimensional.
\end{theorem}

\begin{proposition}[Reconstrucción intrínseca y autosuficiencia probatoria]
\label{prop:frontal-autosuficiencia-probatoria}
Cada resultado \(y\) publicado en este volumen se reconstruye mediante el
recibo matemático local
\[
 \mathscr R_y=
 \bigl(D_y,A_y,\tau_y,U_y,C_y,X_y^{\rm enr},
       y_{\rm HMT},\operatorname{Rec}_y,\mathcal F_y\bigr),
\]
donde \(D_y\) es el dominio; \(A_y\) especifica las hojas y operaciones de
APP; \(\tau_y\) el régimen y la orientación TRIT; \(U_y\) la composición TPK
con dominio, codominio y acción; \(C_y\) los datos conservados de residuo,
cociente, acarreo, hoja, ruta, frontera, memoria y supervivencia;
\(X_y^{\rm enr}\) el estado enriquecido de residencia; \(y_{\rm HMT}\) la
salida anterior a toda identificación externa; \(\operatorname{Rec}_y\) el
reconocimiento convencional posterior; y \(\mathcal F_y\) un falsador de la
flecha sustantiva. La sucesión
\[
 \mathcal D_9\longrightarrow\Omega_{\rm seed}
 \longrightarrow\mathcal U_6^{\rm cat}
 \longrightarrow\mathcal W_6
 \longrightarrow X_\infty^{\rm enr}
 \longrightarrow\mathcal C_{\rm cont}^{\rm disc}
\]
queda definida y demostrada por las ecuaciones, teoremas y recibos impresos
en el propio volumen. Los programas reproducibles anexos vuelven a evaluar
esas identidades; no aportan una premisa ausente ni un valor objetivo.

Los falsadores primarios son: el fallo de la identidad de fibras de
\eqref{eq:frontal-censo-semillas}; la pérdida de alguno de los campos de
\(C_y\) bajo \(U_y\); la incompatibilidad entre extensión y truncamiento; el
vacío del límite inverso; la falta de contracción o de sobreyectividad interna
de los cilindros racionales; la ruptura de naturalidad en la rama
excepcional; y, para una publicación terminal, el incumplimiento de la
hipótesis que define su dominio exacto. Por consiguiente, la validez de una
salida se decide dentro de su dominio y por su falsador explícito, no por una
analogía, una nomenclatura posterior ni una evaluación metrológica.
\end{proposition}

\paragraph{Unicidad del sistema generador y tipado de sus proyecciones.}
Existe una sola APP, seguida por TRIT y por el TPK con estado enriquecido.
APP produce las dos hojas y sus registros; TRIT los orienta; el TPK compone
las rutas y acumula hoja, oposición, residuo, cociente, acarreo, frontera y
memoria en \(\mathcal X_{\rm TPK}^{\rm enr}\). La carta firmada es una
proyección de ese mismo estado terminal:
\[
 R_{\rm sgn}:\mathcal X_{\rm TPK}^{\rm term,enr}
 \longrightarrow\mathcal U_{12}^{\rm sgn},\qquad
 U_{12}^{\rm sgn}=R_{\rm sgn}(x_{\rm term}),\qquad
 K=\frac14\mathcal H_{12}U_{12}^{\rm sgn},\qquad
 \mathcal H_{12}^2=4I_{12}.
\]
Los dos últimos morfismos son enteros y reversibles sobre su imagen. No se
postula un emisor APP directo hacia \(U_{12}^{\rm sgn}\), porque tal flecha
saltaría TRIT y TPK: el dominio correcto de \(R_{\rm sgn}\) es el estado
enriquecido ya generado. Las sombras que olvidan el registro genealógico o la memoria son
proyecciones posteriores, no perfiles matemáticos alternativos. Esta única
composición excluye de sus entradas
\(U,K,D_3K,D_4K,Q,\alpha\) y toda expansión objetivo.

\paragraph{Número y partícula.}
Un número HMT es una sección compatible
\(
 \boldsymbol x=(x_n)_{n\geq1}\in
 \varprojlim_n X_n^{\mathrm{enr}}
\)
con hoja, orientación, residuo, cociente, acarreo, ruta, frontera y memoria.
Su valor arquimediano es la intersección única
\(
 \operatorname{val}(\boldsymbol x)=\bigcap_n I_n(x_n)
\);
el cociente por igualdad de valor publica el real y su fibra conserva las
genealogías que lo realizan. Una partícula HMT--MD es una clase de secciones
persistentes de la extensión enriquecida: el estado conserva la identidad,
la ruta registra su historia observable y los lectores posteriores publican
masa, carga, espín y estadística en sus codominios físicos.

\paragraph{Tipado del corredor excepcional.}
Las lecturas de Hadamard y de Paley parten del mismo estado terminal, pero no
se serializan entre sí. Con
\(X_\infty^{\rm enr}:=\varprojlim_nX_n^{\rm enr}\), se tienen mapas distintos
\begin{equation}
 \operatorname{Pub}_\alpha:X_\infty^{\rm enr}\longrightarrow\mathbb R_{>0},
 \quad x_\infty\longmapsto\alpha_{\rm HMT},
 \qquad
 \operatorname{Inc}_\alpha:X_\infty^{\rm enr}\longrightarrow A_2^{12},
 \quad x_\infty\longmapsto
 v_\alpha=4\rho_1+\rho_2+\cdots+\rho_{12}.
 \label{eq:frontal-dos-publicaciones-alpha}
\end{equation}
Comparten genealogía, pero no se postula igualdad entre sus salidas. La rama
excepcional se bifurca después de \(C_W\):
\begin{equation}
 \begin{aligned}
 C_P&\longrightarrow A_W=C_P\bmod3\longrightarrow
 C_W=[12,6,6]_3\\
 &\longrightarrow
 \begin{cases}
 W_{12}\xrightarrow{\ \widehat\iota_{D,\ell}\ }
 W_{24}\subset\mathcal G_{24},\\[2pt]
 N(C_W)\xrightarrow{\ \operatorname{Nbr}_{v_\alpha}\ }
 N_\alpha\cong\Lambda_{24}\longrightarrow
 V_{\Lambda_{24}}\longrightarrow V^\natural .
 \end{cases}
 \end{aligned}
 \label{eq:frontal-cadena-excepcional}
\end{equation}
Aquí
\begin{equation}
 \mathcal O_{24}:=\{O\subset[24]:|O|=8,\ \mathbf1_O\in\mathcal G_{24}\},
 \qquad
 \mathcal H(D):=\{O\cap D:O\in\mathcal O_{24},\ |O\cap D|=6\}.
 \label{eq:frontal-octadas-hexadas-dodecada}
\end{equation}
Fijado \(\ell:W_{12}\simeq(D,\mathcal H(D))\),
\(\widehat\iota_{D,\ell}(H)=O_H\), donde
\(O_H\cap D=\ell(H)\). Tal octada es única: dos elecciones contendrían el
mismo subconjunto de cinco puntos, contradiciendo \(S(5,8,24)\). La rama
ternaria pega \(N(A_2^{12})=N(C_W)\); el vecino de índice tres determinado
por \(v_\alpha\) produce Leech. Finalmente,
\begin{equation}
 \operatorname{Aut}(V^\natural)\cong\mathbb M,\qquad
 a:\mathbb M\times V^\natural\longrightarrow V^\natural,\qquad
 T_g(\tau)=\operatorname{Tr}_{V^\natural}(gq^{L_0-1}).
 \label{eq:frontal-accion-monstruo}
\end{equation}
Las identidades \(v_\alpha^2=54\) y \([q]J=196884\) son evaluaciones
escalares; ninguna define la acción \(a\).

\paragraph{Realización paralela de la dualidad T.}
La dualidad T no continúa la rama Moonshine: actúa en su propio dominio
\begin{equation}
 \mathscr D_{\rm dual}=\mathbb Z^2\times\mathbb R_{>0},\qquad
 \mathsf T(m,w,R_0)=(w,m,R_0^{-1}),\qquad
 \mathsf T^2=\operatorname{id},\qquad
 E_{R_0}(m,w)=E_{R_0^{-1}}(w,m).
 \label{eq:frontal-dualidad-t}
\end{equation}

\begin{proof}
El censo de \eqref{eq:frontal-censo-semillas} construye el catálogo finito y
selecciona, sin valores objetivo, los tres bloques iniciales \(w_6\). El
registro orientado de transiciones del estado TPK completo explicita los doce
bloques posteriores a la semilla de las tres biografías publicadas. Para cada uno de los dos
intervalos de transición, las seis filas independientes extraídas de esas
tres biografías forman matrices
\(X_\ell,Y_\ell\in M_6(\mathbb F_3)\),
\(\ell\in\{0,1\}\), cuyas filas son, respectivamente, los estados de salida
y de llegada del registro. El cálculo exacto sobre este registro prueba
\(\det X_0=1\), \(\det X_1=-1\) y
\(X_\ell L_\ell=Y_\ell\); por tanto reconstruye unívocamente
\(L_\ell=X_\ell^{-1}Y_\ell\). El
calendario \((L_0,L_0,L_1,L_1)\) es el único de los dieciséis calendarios
binarios que conserva simultáneamente las tres biografías seleccionadas por
los invariantes de fibra; las ciento cuarenta y cuatro perturbaciones
unitarias destruyen ese reconocimiento triple. Por tanto, los transportes,
los doce bloques y su calendario quedan deducidos en la composición única
APP--TRIT--TPK: APP produce las hojas y el registro de partida, TRIT fija la
orientación de cada transición y TPK prolonga el registro con memoria. El
censo de semillas y el estado enriquecido son etapas sucesivas del mismo
generador, no perfiles rivales ni teorías diferentes.

La extensión no escindida
\[
0\longrightarrow C_{12}\longrightarrow C_{108}
\longrightarrow C_9\longrightarrow0
\]
separa retorno de fase y retorno de estado. Los truncamientos eliminan la
última actualización y conservan el prefijo; el lema de König produce una
rama infinita compatible. Para cada precisión \(N\), la recurrencia y la
holonomía generan la sección y sus caracteres, y determinan el cilindro que
prolonga el registro genealógico anterior. La familia de prefijos satisface
\[
\rho_{N+1,N}(p_{N+1})=p_N,
\qquad
p_\infty\in\varprojlim_N\{0,\ldots,9\}^{N}.
\]
Mil, diez mil o un millón de cifras son proyecciones finitas de este único
objeto infinito.

Las horquillas racionales anidadas
\(
 p_N^-<\chi<p_N^+
\)
se calculan después de producir el carácter \(\chi\) y satisfacen
\(
 [p_{N+1}^-,p_{N+1}^+]\subset[p_N^-,p_N^+]
\)
con diámetro tendente a cero. Certifican la salida generada a cualquier
precisión finita; no seleccionan la órbita, el transporte ni el carácter.

Las tres vueltas enriquecidas \(\gamma_{-},\gamma_0,\gamma_+\) de nueve
emisiones se distinguen por el acarreo
\(\operatorname{Carr}(\gamma_\rho)=\rho\kappa_{\rm Carr}\) y admiten las
tres prolongaciones después de cualquier prefijo. El límite de sus rutas
registradas satisface
\[
 \operatorname{dig}_{\rm Carr}^{\omega}:
 \Omega_\Gamma^{\rm cal}\xrightarrow{\sim}
 \{-1,0,+1\}^{\mathbb N}.
\]
Tras agrupar seis retornos, la aplicación
\[
 \beta=\operatorname{blk}_6\circ
 (\rho\mapsto\rho+1)^{\mathbb N}\circ
 \operatorname{dig}_{\rm Carr}^{\omega}:
 \Omega_\Gamma^{\rm cal}\xrightarrow{\sim}
 (\mathbb F_3^6)^{\mathbb N}
\]
es un homeomorfismo; su inversa concatena las vueltas prescritas. Las formas
equilibradas normalizadas por
\[
 a_k+c_k=\rho_k+3c_{k+1},\qquad
 \rho_k\in\{-1,0,+1\},
\]
producen primero el anillo ordenado \(\mathbb Z_{\rm HMT}\), después su
cuerpo de fracciones \(\mathbb Q_{\rm HMT}\) y finalmente
\[
 \mathbb R_{\rm HMT}:=
 \operatorname{Cau}(\mathbb Q_{\rm HMT})/{\asymp}.
\]
Si \(\beta(\xi)=(b_q)\) y
\(\nu(b)=\sum_{i=1}^{6}b_i3^{6-i}\), la sucesión
\[
 q_n=m+\sum_{q=0}^{n-1}\frac{\nu(b_q)}{729^{q+1}}
\]
es de Cauchy, y la partición interna recurrente en \(729\) subintervalos
prueba la sobreyectividad
\[
 P_{\rm ar}:\mathbb Z_{\rm HMT}\times\Omega_\Gamma^{\rm cal}
 \twoheadrightarrow\mathbb R_{\rm HMT},\qquad
 (\mathbb Z_{\rm HMT}\times\Omega_\Gamma^{\rm cal})/{\sim_{\rm ar}}
 \cong\mathbb R_{\rm HMT}.
\]
Sólo después, la unicidad del cuerpo ordenado completo arquimediano publica
\(\mathbb R_{\rm HMT}\cong\mathbb R\). La dirección interior subdivide
\([0_{\rm HMT},1_{\rm HMT})\); la exterior concatena el prefijo que codifica
\(m\in\mathbb Z_{\rm HMT}\). Así se cubre la recta sin usar previamente
\(\mathbb R\), parte entera ni expansión decimal objetivo.
Los subespacios de supervivencia de caracteres particulares individualizan
secciones dentro de esta capacidad total. La cobertura del continuo procede
de la completación global; la recurrencia nonádica genera
\(\pi_{\mathrm{HMT}},e_{\mathrm{HMT}},\varphi_{\mathrm{HMT}},
\alpha_{\mathrm{HMT}}\), y los filtros tipados certifican sus biografías
dentro de esa cobertura.

Las cinco acciones espectral, solenoidal, gauge, cohomológica y
determinantal se calculan sobre una misma emisión enriquecida y comparten
soporte, ruta, registro genealógico y truncamientos. Su naturalidad conjunta las conserva
en el límite. La dicotomía perfecta clasifica cada publicación analítica
como numerable o de cardinal continuo. El código generado \(h\) selecciona
el universo genealógico canónico \(M_h=L[h]\). La construcción de los
reales HMT identifica cada número por valor, cilindro, ruta y memoria; el
buen orden inducido por esa genealogía prueba la fórmula clásica
\[
M_h\models2^{\aleph_0}=\aleph_1.
\]
Como \(M_h\) es precisamente el codominio generado por la construcción del
continuo de este teorema, la conclusión terminal del sistema es
\[
 \boxed{2^{\aleph_0}=\aleph_1.}
\]
La independencia en el lenguaje desnudo de ZFC clasifica la pérdida del
selector genealógico \(h\); HMT conserva ese dato y determina el cardinal de
su continuo. La sentencia conservada es la hipótesis clásica del continuo,
con los cardinales ordinarios de \(M_h\), y la prueba detallada establece el
puente entre cilindros compatibles, buen orden genealógico y cardinalidad.
Este cierre cardinal se distingue del corredor lógico fuerte, también
construido desde el continuo genealógico:
\[
\begin{aligned}
 \mathcal C_{\mathrm{cont}}^{\mathrm{disc}}
 &\longrightarrow\mathcal L_0^{\mathrm{log}}
 \xrightarrow{\Phi}\mathcal L_1^{\mathrm{log}}
 \longrightarrow\gamma_F\longrightarrow w(\gamma_F)\\
 &\longrightarrow\mathsf{NoEscape}_{24}
 \longrightarrow\mathsf{AntiDiag}_{369}
 \longrightarrow\mathsf{G\ddot{o}delIII}_{\mathrm{HMT}}\\
 &\longrightarrow\mathsf{CycleCriterion}_{\mathrm{ZFC}},
\end{aligned}
\]
que publica, en su dominio tipado, el criterio
\[
 \boxed{
 \neg\operatorname{Consis}(\mathrm{ZFC})
 \Longleftrightarrow
 \exists C\;\bigl(C\text{ es un ciclo fuera de banda y }|C|>24\bigr).}
\]
Las letras \(\mathcal L_0^{\mathrm{log}},\mathcal L_1^{\mathrm{log}}\)
denotan aquí los objetos del corredor lógico y no las matrices de elevación
ternaria \(L_0,L_1\); la separación de tipos impide que una conclusión
local sobre uno de esos pares debilite la del otro.

Las tres órbitas funcionales del catálogo producen los caracteres HMT de
clausura, propagación y autoescala. Sus leyes internas fuerzan,
respectivamente,
\[
16\arctan(1/5)-4\arctan(1/239)=\pi_{\mathrm{HMT}}=\pi,
\qquad
P_{s+t}=P_sP_t\Rightarrow
P_1=\sum_{n\ge0}\frac1{n!}=e_{\mathrm{HMT}}=e,
\]
\[
F_{\mathrm{av}}\binom1{\varphi_{\mathrm{HMT}}}
=\varphi_{\mathrm{HMT}}\binom1{\varphi_{\mathrm{HMT}}},
\qquad \varphi_{\mathrm{HMT}}=\varphi.
\]
La involución interna intercambia las ramas
\(\pi_{\mathrm{HMT}}/e_{\mathrm{HMT}}\), conserva el eje
\(\varphi_{\mathrm{HMT}}\) y retiene el agregado
\(e_{\mathrm{HMT}}+\pi_{\mathrm{HMT}}\). En la fase TPK enriquecida de la
única composición, la proyección canónica del estado terminal publica su coordenada primitiva
firmada
\(U=R_{\mathrm{sgn}}(x_\infty)\), y la identidad
\(\mathcal H_{12}^2=4I\) reconstruye de manera entera y única el sello
\(K=\mathcal H_{12}U/4\). La recurrencia dodecafásica con acarreo y la
prolongación remota de sus cilindros publican \(\alpha_{\rm HMT}\). De
manera independiente, el estado enriquecido determina la firma
\[
 \mathcal I_9=(2\pi_{\rm HMT},3,4,7,20,21,45,46,120,11),
\]
y el lector graduado asociado construye, sin consultar la raíz, el operador
polinómico exacto
\(E_{21/22}^{(9)}=P_9-\log_{10}(10/9)\). Su única raíz positiva simple
\(\xi_9\) es el testigo finito de orden nueve de la segunda publicación
analítica interna. El mismo transporte graduado del TPK se itera sobre los
grados sucesivos y construye el sistema compatible completo
\(E_{21/22}^{(6+3m)}\); sus cuadrados de truncamiento conmutan con los de la
familia dodecafásica y determinan, a cada profundidad, el mismo cilindro
superviviente. En consecuencia,
\[
 \alpha_A:=\varprojlim_N\rho_N(\alpha_{12}),\qquad
 E_{21/22}:=\varprojlim_mE_{21/22}^{(6+3m)},\qquad
 \alpha_B:=\operatorname*{arg\,zero}_{x\in I_\alpha}E_{21/22}(x),
\]
\[
 \boxed{\alpha_A=\alpha_B=\alpha_{\rm HMT}.}
\]
La igualdad finita
\[
 \operatorname{Tr}_9(\xi_9^{-1})
 =\operatorname{Tr}_9(\alpha_{12}^{-1})=137.035999084
\]
certifica el primer cuadrado visible de esa identidad; no la define, no la
limita a nueve cifras y no selecciona los coeficientes de la prolongación.
La clausura bidireccional funcional asociada es
\[
\alpha T^2-B_9(\alpha)T+\alpha^3=0,
\qquad
T_\pm=\alpha e^{\pm\xi},
\qquad
T_+T_-=\alpha^2,
\qquad
J_\alpha(T)=\frac{\alpha^2}{T}.
\]
Así, la coordenada publicada por la clausura dodecafásica es la media
geométrica conservada por las dos rutas conjugadas. La vía dodecafásica y el
lector analítico mantienen dominios, operadores y falsadores propios, pero
sus sistemas inversos publican una misma sección límite. El comparador
finito y el reconocimiento metrológico son posteriores.

La coordenada incidencial del sector \(\alpha\) fija
\[
v_\alpha=4\rho_1+\rho_2+\cdots+\rho_{12},
\qquad v_\alpha^2=54,
\]
y el vecino de índice tres
\[
N_{\alpha,0}
=\{x\in N(C_W):\langle x,v_\alpha\rangle\equiv0\pmod3\},
\qquad
N_\alpha=N_{\alpha,0}+\mathbb Z\frac{v_\alpha}{3}
\cong\Lambda_{24}.
\]
Esta rama construye exactamente el vecino de Leech y, después, la cadena
\(\Lambda_{24}\to V^\natural\) y la acción de
\(\operatorname{Aut}(V^\natural)\cong\mathbb M\) sobre \(V^\natural\). El
escalar \(\alpha_{\mathrm{HMT}}\) y el vector \(v_\alpha\) son las dos
coordenadas tipadas del mismo sector dodecafásico. Su acoplamiento exacto es
la bisagra de origen común que conduce a las simetrías del Monstruo: la
coordenada escalar se publica en \(\mathbb R_{\rm HMT}\) y la coordenada
incidencial se publica en la rama
\(A_2^{12}\to\Lambda_{24}\to V^\natural\to\mathbb M\), con sus codominios
conservados y su antecedente genealógico compartido explícito.

La Mecánica Dimensional recibe las salidas anteriores y compone, en su orden
constitutivo, la retención de fase, la década \(-34\), la pareja de Planck,
los canales \(90/120\), el electrón, el atlas de masas, la gravedad, el área
y la información. Los capítulos siguientes demuestran cada flecha, publican
su falsador y enlazan su certificado ejecutable. Las identificaciones
convencionales comparecen después como reconocimiento de las magnitudes ya
producidas. Esto completa la cadena de
\eqref{eq:frontal-doble-proyeccion}.
\end{proof}

\section*{Jerarquía de las publicaciones geométricas y dimensionales}
\addcontentsline{toc}{section}{Jerarquía de las publicaciones geométricas y dimensionales}

\begin{proposition}[TRIT como unidad local de los tres regímenes]
La firma trítica conserva dos lectores coordinados. El lector geométrico
publica \(\tau_{\rm g}\in\{+1,0,-1\}\) y
\begin{equation}
 J_{\tau_{\rm g}}^{\,2}=-\tau_{\rm g}I,\qquad
 J_+^{\,2}=-I,\qquad J_0^{\,2}=0,\qquad J_-^{\,2}=+I,
 \label{eq:frontal-trit-tres-regimenes}
\end{equation}
y distingue los regímenes elíptico, parabólico e hiperbólico. El lector
temporal activo publica otro signo, \(\varepsilon_{\rm t}\), mediante
\begin{equation}
 \begin{aligned}
 +1&\mapsto\text{retorno, memoria y pasado},\\
 0&\mapsto\text{umbral y presente},\\
 -1&\mapsto\text{apertura bidireccional y futuro}.
 \end{aligned}
 \label{eq:frontal-trit-tiempo}
\end{equation}
No se postula \(\tau_{\rm g}=\varepsilon_{\rm t}\). Estas dos salidas no
identifican signo de curvatura y orientación temporal: son proyecciones
distintas de una misma unidad local. El TPK transporta la
unidad completa en
\begin{equation}
 \mathcal F_q=
 \mathcal O_q\times\mathcal H_q\times\mathcal C_q\times
 \mathcal S_q\times\mathcal B_q\times\Lambda_q
 \label{eq:frontal-fibra-tpk}
\end{equation}
y conserva hoja y orientación, residuo y cociente, acarreo, firma, cilindro,
frontera y registro. En particular, la rama elíptica posee el realizador
finito
\begin{equation}
 A_W^2=2I_6=-I_6\quad\text{en }M_6(\mathbb F_3),\qquad
 i_{\rm HMT}:=\chi_i(A_W),\qquad i_{\rm HMT}^{\,2}=-1,
 \label{eq:frontal-raiz-interna-menos-uno}
\end{equation}
donde \(\chi_i:\langle A_W\rangle\to\mathbb C\) es la carta orientada.
La raíz de \(-1\), el plano complejo y el sistema de Euler son publicaciones
de un endomorfismo construido, no entradas del generador.
\end{proposition}

\begin{proof}
El selector trítico incorpora uno de los tres generadores de
\eqref{eq:frontal-trit-tres-regimenes} sin borrar las hojas suma y producto.
El transporte cardinal conserva su orientación y la actualización conserva
los seis factores de \eqref{eq:frontal-fibra-tpk}. Por tanto, la proyección
real puede condensar los tres regímenes, mientras el estado enriquecido
reconstruye cuál de ellos y qué memoria produjeron el valor visible. Las
realizaciones trigonométrica, nilpotente e hiperbólica desarrolladas en el
cuerpo son las imágenes respectivas de este único tipado local.
\end{proof}

\begin{theorem}[Cadena de acción, gravedad y materia]
Después de publicar correlacionadamente
\((\pi_{\rm HMT},\varphi_{\rm HMT},e_{\rm HMT},\alpha_{\rm HMT})\), la rama
determinantal fija la década
\begin{equation}
 \Sigma_H=\varphi_{\rm HMT}\alpha_{\rm HMT}^{16},\qquad
 d_H:=\lfloor\log_{10}\Sigma_H\rfloor=-34.
 \label{eq:frontal-decada-accion}
\end{equation}
La mantisa procede del lector regional distinto
\begin{align}
 H_5&=\frac12\sqrt{\frac{1000\alpha_{\rm HMT}}{\varphi_{\rm HMT}}}
 -\alpha_{\rm HMT}+\frac9{16}\alpha_{\rm HMT}^2
 -\frac59\alpha_{\rm HMT}^3+\frac7{48}\alpha_{\rm HMT}^4
 -\frac1{54}\alpha_{\rm HMT}^5,\notag\\
 D_A&=\exp\!\left(-\frac{100\pi_{\rm HMT}}9\alpha_{\rm HMT}\right),
 \qquad
 \mathscr C_\pi=169\alpha_{\rm HMT}^6+
 \frac{D_A\alpha_{\rm HMT}^7}{1-D_A\alpha_{\rm HMT}},\notag\\
 \eta_{\rm pre}&:=H_5,\qquad
 \eta_{\rm ret}:=H_5-\frac{90}{\pi_{\rm HMT}}\mathscr C_\pi>0.
 \label{eq:frontal-mantisas-accion}
\end{align}
Para \(a\in\{\mathrm{pre},\mathrm{ret}\}\), las dos secciones de acción son
\begin{equation}
 \hbar_a=\eta_a10^{d_H}\mathcal U_S
 =\eta_a10^{-34}\mathcal U_S,\qquad
 h_a=2\pi_{\rm HMT}\hbar_a.
 \label{eq:frontal-secciones-accion}
\end{equation}
\(\Sigma_H\) fija sólo la década; no se presenta como generador de la
mantisa.

Fijados \(t_0>0\), \(c_{\rm int}>0\) y \(\ell_0=c_{\rm int}t_0\), CT108
publica
\begin{equation}
 \omega_{108}=\frac{2\pi_{\rm HMT}}{108t_0},\qquad
 R_{108}=\frac{c_{\rm int}}{\omega_{108}}
 =\frac{54}{\pi_{\rm HMT}}\ell_0,\qquad
 m_{108,a}=\frac{\hbar_a\omega_{108}}{c_{\rm int}^2}.
 \label{eq:frontal-reloj-gravedad}
\end{equation}
La familia gravitatoria tipada es
\begin{equation}
 G_a(\vartheta)=\vartheta\frac{c_{\rm int}^5t_0^2}{\hbar_a}.
 \label{eq:frontal-familia-gravedad}
\end{equation}
En la convención reducida
\(r_{g,a}=G_a m_{108,a}/c_{\rm int}^2\), el cierre
\(r_{g,a}(\vartheta)=R_{108}\) posee la solución positiva única
\begin{equation}
 \vartheta_{108}=\left(\frac{54}{\pi_{\rm HMT}}\right)^2,\qquad
 G_a:=G_a(\vartheta_{108}),
 \label{eq:frontal-selector-gravedad}
\end{equation}
y sólo entonces publica
\begin{equation}
 \ell_P:=\sqrt{\frac{G_a\hbar_a}{c_{\rm int}^{\,3}}}
 =\sqrt{\vartheta_{108}}\,\ell_0
 =R_{108}=\bar\lambda_{C,a}=r_{g,a}.
 \label{eq:frontal-longitud-planck}
\end{equation}
Además, \(G_a\hbar_a=\vartheta_{108}c_{\rm int}^5t_0^2\), por lo que
\(\ell_P\) es independiente de la orientación \(a\). La gravedad tensorial
\[
 \mathbb G_{{\rm TT},a}(n)
 =G_a\Lambda(n)\Gamma_5P_5:
 V_{12}\longrightarrow\mathcal H_{\rm TT}(n)
\]
es posterior al escalar. La acción Palatini--Cartan--Holst recibe después
\(G_a\) y la salida independiente \(\gamma_{\rm HMT}\); el funcional
hexada--octada transporta esta última al espectro de área y a la información
de frontera.

En la rama de materia, la ley completa actúa sobre \(81\) celdas, \(56\)
familias, \(13\) multisecciones y \(324\) rutas, con sector nulo separado. Los
lectores \(K_D\) y \(K_\Omega\) poseen, respectivamente, \(14\) y \(9\)
perfiles, forman \(40\) pares y coinciden en \(18\) rutas únicamente en
\((80,54,6)\). La fibra central \((5,5)\) produce el electrón beta. Sobre las
fibras restantes, la ley publica operadores positivos, medidas espectrales y
multiplicidades; el reconocimiento exterior se abre después mediante el
inventario exhaustivo
\begin{equation}
 \mathscr I_{\rm ext}
 =\mathscr I_m^{471}\coprod\mathscr I_\Gamma^{384},
 \qquad |\mathscr I_{\rm ext}|=855.
 \label{eq:frontal-inventario-masas-anchuras}
\end{equation}
CKM transporta los marcos espectrales quark y PMNS los marcos cargado y neutro
de rango tres; sus dominios, operadores y falsadores permanecen distintos.
Ningún valor del inventario selecciona una ruta, un coeficiente, un autovalor
o un ángulo.
\end{theorem}

\begin{proof}
La forma normal de Smith produce
\(\operatorname{diag}(1,2,2,4)\) y conúcleo de orden \(16\); las cotas de
\(\Sigma_H\) fijan \(d_H=-34\), mientras
\eqref{eq:frontal-mantisas-accion}--\eqref{eq:frontal-secciones-accion}
publican \(\hbar_a\) y después \(h_a\). La sustitución de
\(\vartheta_{108}\) en \(G_a(\vartheta)\) prueba
\eqref{eq:frontal-longitud-planck}. La incidencia de los conjuntos
\(H_{\rm act}\times\binom{H_{\rm act}}2\) y
\(\{1,\ldots,8\}\times\binom{H_{\rm act}}2\) produce \(90\) y \(120\).
Los doce sectores dodecafásicos y la única costura orientada de retorno
producen, mediante el lector lineal, el peso \(12{:}{-}1\) de la transición
hexada--octada. El coeficiente de \((1-q)^{-1}\), igual a \(1/24\), y la
minimalidad diofántica verifican posteriormente ese peso. El censo finito del atlas y de sus dos lectores
prueba los cardinales de la rama de masa; la inversión celular tiene un único
punto fijo, \((5,5)\), que reduce canónicamente la fibra electrónica a rango
uno. El inventario externo se compone sólo después de congelar fibra,
operador, proyector, representación, esquema y escala, de modo que una
permutación de sus magnitudes deja invariante todo el constructor interno.
\end{proof}
% END HMT-MD-FRONTAL-MAXIMAL-SIGNATURES-20260904

\fi
